\Article{247}
\textsc{Satz.}\quad\emph{Das Produkt zweier gleichartiger Perioden lässt sich unabhängig von der Zahl $p$ durch Addition gleichartiger Perioden und gegebener Zahlen bilden.}

(Perioden heißen gleichartig, wenn sie gleich viele Terme enthalten oder die Zahl $e$ dieselbe ist.)

Beispiel. Es sei $n=7$. Das Produkt der Perioden $(1)+(6)$ und $(2)+(5)$ ist wegen $(a)\times(b)=(a+b)$ gleich $(3)+(6)+(8)+(11)$, besteht also aus den Perioden $(3)+(4)$ und $(1)+(6)$.

\emph{Beweis.}\quad Es sei $\frac{\lambda}{e}=f$, und die gegebenen Perioden seien $(e\ast\alpha)$ und $(e\ast\beta)$ oder die Summen
\[
\begin{aligned}
&(\rho^\alpha)+(\rho^{\alpha+e})+(\rho^{\alpha+2e})+\ldots+(\rho^{\alpha+(f-1)e})\ \dotsfive\ P\\
&(\rho^\beta)+(\rho^{\beta+e})+(\rho^{\beta+2e})+\ldots+(\rho^{\beta+(f-1)e})\ \dotsfive\ Q
\end{aligned}
\]

Das Produkt $PQ$ besteht aus $f^2$ Termen. Diese ordne man folgendermaßen an. Man bilde $f$ Zeilen, deren jede aus $f$ Termen besteht. Die erste enthalte das Produkt von $P$ mit $(\rho^\beta)$, die zweite das Produkt $P\pdot(\rho^{\beta+e})$ usw. usw. In der ersten Zeile nehme das aus dem Term $(\rho^\alpha)$ entstehende Produkt die erste Stelle ein, das aus $(\rho^{\alpha+e})$ entstehende die zweite, und so weiter der Reihe nach. In der zweiten Zeile weise man die erste Stelle dem aus $(\rho^{\alpha+e})$ entstehenden Produkt zu, die zweite dem aus $(\rho^{\alpha+2e})$ entstehenden usw. und schließlich die letzte dem aus $(\rho^\alpha)$ entstehenden. Die dritte beginne mit dem aus $(\rho^{\alpha+2e})$ entstehenden Produkt, und so fahre man fort; auf das Produkt aus dem letzten Term folgen die Produkte aus dem ersten und zweiten usw. usw. Bezeichnet man also die aufeinanderfolgenden Terme der Periode $P$ mit\vadjust{\penalty10000} $1,2,3\ldots z$ und die der Periode $Q$ mit $\mathrm{I},\mathrm{II},\mathrm{III},\dotsfour Z$, so ordne man die Teile von $PQ$ folgendermaßen an:
\[
\begin{aligned}
&1\pdot\mathrm{I}+2\pdot\mathrm{I}+3\pdot\mathrm{I}+4\pdot\mathrm{I}+\dotstwo+z\pdot\mathrm{I}\\
&2\pdot\mathrm{II}+3\pdot\mathrm{II}+4\pdot\mathrm{II}+\dotsfour+1\pdot\mathrm{II}\\
&3\pdot\mathrm{III}+4\pdot\mathrm{III}+\ldots+1\pdot\mathrm{III}+2\pdot\mathrm{III}\\
&\qquad\text{usw.}\quad\text{usw.}
\end{aligned}
\]

Dann fasse man alle Terme, die in den einzelnen Zeilen dieselbe Stelle einnehmen, zu $f$ Gruppen zusammen. Ich behaupte:
\SourceLeaf{216}
\par\noindent $1^\circ$ Ist irgendein Term $\equiv1$, so sind alle übrigen Terme derselben Gruppe ebenfalls $\equiv1$

\noindent $2^\circ$ Jede Gruppe, in der kein Term $\equiv1$ ist, bildet eine Periode. --- Sobald dies bewiesen ist, haben wir unser Ziel offensichtlich erreicht.

Die allgemeine Form einer solchen Gruppe ist
\[
\begin{gathered}
(\rho^{\alpha+ke}+\rho^\beta),\quad(\rho^{\alpha+(k+1)e}+\rho^{\beta+e}),\quad(\rho^{\alpha+(k+2)e}+\rho^{\beta+2e}),\\
\ldots\ (\rho^{\alpha+(k+f-1)e}+\rho^{\beta+(f-1)e})
\end{gathered}
\]
Denn für $\rho^{\alpha+(k-1)e}$ kann auch $\rho^{\alpha+(k+f-1)e}$ geschrieben werden, weil $ef=\lambda$ und $\rho^\lambda\equiv1\;(\mathrm{mod.}\,n)$ gelten, und ebenso für die vorhergehenden Terme. Man setze $\rho^{\alpha+ke}+\rho^\beta\equiv\rho^\chi\;(\mathrm{mod.}\,n)$, was erlaubt ist, sofern nicht $\rho^{\alpha+ke}+\rho^\beta$ durch $n$ teilbar ist\begingroup\renewcommand{\thefootnote}{*)}\footnote{Die Aussage muss etwas anders formuliert werden, wenn $n$ allgemein eine Primzahlpotenz bezeichnet; ist es dagegen eine Primzahl, so muss nichts geändert werden.}\endgroup. Dann lässt sich die Gruppe so darstellen: $(\rho^\chi)$, $(\rho^{\chi+e})$, $(\rho^{\chi+2e})\ldots(\rho^{\chi+(f-1)e})$. Sie ist offensichtlich die Periode $(e\ast\chi)$. Ist dagegen $\rho^{\alpha+ke}+\rho^\beta$ durch $n$ teilbar, so sind alle Terme der Gruppe $\equiv(0)$, also $\equiv1$. Q.\,E.\,D.

\emph{Anmerkung.} Dieser Beweis zeigt zugleich ein sehr einfaches Verfahren, um das Produkt zu entwickeln. Weiter unten geben wir ein anderes, das zwar diesen Vorzug nicht hat, aber wegen seiner Einfachheit nicht unbeachtlich erscheint.

\Needspace{15\baselineskip}\Article{248}
Alle kleineren Perioden, die eine größere Periode bilden, nennen wir ein Periodensystem. So erhalten die Perioden
\[
(ef\ast\alpha),\quad(ef\ast f+\alpha),\quad(ef\ast2f+\alpha)\ \dotsfour\ (ef\ast(e-1)f+\alpha)
\]
aus denen die Periode $(f\ast\alpha)$ besteht, diesen Namen. Das System ist \emph{ordnungsgemäß angeordnet}, wenn die Zahlen hinter dem Zeichen $\ast$, hier $\alpha$, $f+\alpha$, $2f+\alpha$, eine arithmetische Folge mit Differenz $f$ bilden. Schließlich heißen die Systeme \emph{gleichartig}, wenn sowohl die kleineren als auch die größeren Perioden gleichartig sind.

\textsc{Satz.} \emph{Werden die Perioden zweier gleichartiger, ordnungsgemäß angeordneter Systeme miteinander multipliziert, nämlich die erste mit der ersten, die zweite mit der zweiten, die dritte mit der dritten usw., so lässt sich die Summe aller Produkte aus gegebenen Zahlen und Perioden bilden, die der größeren Periode gleichartig sind.}

\emph{Beweis.} Die Systeme seien
\[
\begin{gathered}
(ef\ast\alpha),\quad(ef\ast\alpha+f),\quad(ef\ast\alpha+2f)\ldots\\
(ef\ast\beta),\quad(ef\ast\beta+f),\quad(ef\ast\beta+2f)\ldots
\end{gathered}
\]

\SourceLeaf{217}
Die Produkte der einzelnen Perioden des ersten Systems mit den entsprechenden Perioden des zweiten bestehen nach dem vorigen Abschnitt aus ganzen Zahlen und gleichartigen Perioden. Betrachtet man die Entstehung dieser Perioden mit etwas Aufmerksamkeit, so erkennt man: Besteht $(ef\ast\alpha)\times(ef\ast\beta)$ aus einer ganzen Zahl $N$ und den Perioden $(ef\ast A)$, $(ef\ast B)$, $(ef\ast C)$ usw., so ergeben sich die Produkte
\[
\begin{aligned}
&(ef\ast\alpha+f)\times(ef\ast\beta+f)\text{ aus }N\text{ und Perioden}\\
&\hspace{2em}(ef\ast A+f),\quad(ef\ast B+f),\quad(ef\ast C+f)\text{ usw.}\\
&(ef\ast\alpha+2f)\times(ef\ast\beta+2f)\text{ aus }N\text{ und Perioden}\\
&\hspace{2em}(ef\ast A+2f),\quad(ef\ast B+2f),\quad(ef\ast C+2f)\text{ usw.}
\end{aligned}
\]
und allgemein
\[
\begin{aligned}
&(ef\ast\alpha+\mu f)\times(ef\ast\beta+\mu f)\text{ aus }N\text{ und Perioden}\\
&\hspace{2em}(ef\ast A+\mu f),\quad(ef\ast B+\mu f),\quad(ef\ast C+\mu f)\text{ usw.}
\end{aligned}
\]
Daraus folgt unmittelbar, dass die Summe aller Perioden
\[
eN+(f\ast A)+(f\ast B)+(f\ast C)\text{ usw.}\qquad\text{Q.\,E.\,D.}
\]
ist. Auch dieser Beweis liefert ein Verfahren, um jene Summe zu finden.

\Article{249}
Dieser Satz lässt sich leicht noch verallgemeinern: Hat man beliebig viele gleichartige, ordnungsgemäß angeordnete Systeme und bildet die Produkte aller ersten Perioden, aller zweiten usw., so besteht die Summe all dieser Produkte aus Zahlen und größeren Perioden. Nimmt man alle diese Systeme als gleich an, so besteht die Summe der Potenzen aller Perioden bei beliebigem Exponenten aus Zahlen und Perioden, die der größeren gleichartig sind. Schon daraus wird das Ziel dieser Untersuchungen klar. Es sei $\lambda=efgh\ldots$. Man zerlege alle primitiven Wurzeln in $e$ Perioden $A,A',A''$ usw., jede dieser Perioden wiederum in $f$: $B,B',B''$ usw., und jede davon in $g$: $C,C',C''$ usw. Die Summe aller Perioden ist gegeben, nämlich $\equiv-1$. Nach dem gerade Dargelegten sind aber auch
\[
\begin{gathered}
(A)^2+(A')^2+(A'')^2+(A''')^2+\text{usw.}\\
(A)^3+(A')^3+(A'')^3+(A''')^3+\text{usw.}\\
\text{usw. usw.}
\end{gathered}
\]
gegeben.

Daraus lässt sich nach Abschnitt 244 eine Kongruenz des Grades $e$ finden, deren Wurzeln $A,A',A''$ usw. sind. Setzt man diese nun als bekannt voraus, so zerlege man jede Periode in kleinere:
\SourceLeaf{218}
\[
\begin{array}{rclll}
A&\text{in}&B,&B',&B''\ldots\\
A'&\text{in}&B^{(n)},&B^{(n+1)},&B^{(n+2)}\ldots\\
&\text{usw.}
\end{array}
\]
Also ist $B+B'+B''+\ldots\equiv A$ gegeben. Nun bestehen aber
\[
\begin{gathered}
(B)^2+(B')^2+(B'')^2+\ldots\\
(B)^3+(B')^3+(B'')^3+\ldots\\
\text{usw.}
\end{gathered}
\]
aus Einheiten und den Perioden $A,A',A''$ usw. Daher sind $B,B',B''$ usw. durch eine Kongruenz des Grades $f$ gegeben, aus der sie gefunden werden können. Ebenso lassen sich die Perioden bestimmen, aus denen $A',A''$ usw. bestehen. Daraus wird jeder erkennen, dass sich jede Periode auf genau dieselbe Weise in kleinere unterteilen lässt, bis man zu den Wurzeln selbst gelangt.

\Article{250}
Bei der Anwendung dieser Regeln tritt aber eine Schwierigkeit auf, die wir beseitigen müssen. Da jede Kongruenz mehrere Wurzeln hat, müssen wir sehen, welches Zeichen jeder zuzuweisen ist, damit sie sich richtig voneinander unterscheiden lassen. Die Bezeichnung der Perioden hängt von den Zahlen $r,\rho$ ab, die beliebig gewählt werden können. Daher muss auch der Bezeichnung etwas Willkürliches anhaften. Die Zahl $\rho$ muss allerdings schon zu Beginn festgelegt werden. Das Wesen unserer Methode besteht vor allem darin, aus größeren Perioden kleinere abzuleiten. Ohne die richtige Anordnung der Perioden, die wir durch die \emph{Zeichen} erreicht haben, ist das aber nicht möglich. Wir müssen daher dafür sorgen, dass alle Perioden, sobald sie gefunden sind, durch ihre Zeichen unterschieden werden.

Die Periode $A$ sei mit $(e\ast\alpha)$ bezeichnet und in $f$ Perioden $B,C,D$ usw. zerlegt, die bezeichnet werden sollen. Jede ist offenbar in einer solchen Form $(ef\ast ke+\alpha)$ enthalten. Ich behaupte aber, dass man für eine von ihnen, $B$, die Zahl $k$ beliebig wählen und daraus die Anordnung der übrigen ableiten kann.

Es sei $R$ eine beliebige primitive Wurzel der Kongruenz $x^n\equiv1$, und $B$ bestehe aus den Termen $R^\mu+R^\nu+\text{usw.}$. Es sei $\frac1\mu\rho^{ke+\alpha}\equiv\frac1\upsilon\;(\mathrm{mod.}\,n)$. Da der Wert von $r$ beliebig ist, sofern nur $A$ das Zeichen $(e\ast\alpha)$ erhält, was offensichtlich von selbst geschieht, setze man $r\equiv R^\upsilon\;(\mathrm{mod.}\,p)$. Dann ist der erste Term von $B$ gleich $r^{\rho^{ke+\alpha}}$, und $B$ darf mit
\SourceLeaf{219}
$(ef\ast ke+\alpha)$ bezeichnet werden. Hätten wir anstelle von $R^\mu$ den Term $R^\nu$ betrachtet, so hätten wir einen anderen Wert von $r$ erhalten. Man erkennt aber leicht, dass für jede beliebige Wurzel $\rho$ die Wurzel $r$ $\frac\lambda{ef}$ verschiedene Werte haben kann.

\Article{251}
Wir wollen nun sehen, wie man nach der Bezeichnung einer Periode die übrigen durch ihre Zeichen unterscheidet. Dazu müssen wir aber eine andere Methode suchen, um die übrigen Perioden zu finden. Denn soweit diese wie $A$ selbst Wurzeln einer Kongruenz sind, ist unter ihnen keine Reihenfolge erkennbar. Wir nehmen an, $A$ sei mit $(ef\ast0)$ bezeichnet. Aus dem Vorhergehenden folgt dann
\[
\begin{aligned}
A^2&\text{ hat die Form }M+N(ef\ast0)+O(ef\ast1)+P(ef\ast2)+\ldots\\
A^3&\text{ hat die Form }M'+N'(ef\ast0)+O'(ef\ast1)+\ldots\\
&\text{usw.}\\
A^{ef-1}&\text{ hat die Form }M^\ast+N^\ast(ef\ast0)+O^\ast(ef\ast1)+\ldots
\end{aligned}
\]
Hinzu kommt die Kongruenz
\[
(ef\ast0)+(ef\ast1)+\ldots+(ef\ast ef-1)\equiv-1
\]
Wir haben also $cf-1$ lineare Kongruenzen und ebenso viele unbekannte Größen, die sich daher durch Elimination bestimmen lassen.

\emph{Anmerkung.} Es kann vorkommen, dass die unbekannten Größen durch Ausdrücke der Form $\frac{V}{Wp}$ gegeben sind. Wie dieser Schwierigkeit abzuhelfen ist, zeigen wir weiter unten. Hier wollen wir uns nicht damit aufhalten, weil dieser Fall nur sehr selten auftreten kann.

\Article{252}
Dies möge im Allgemeinen über die Lösung reiner Kongruenzen genügen. An verschiedenen Stellen weiter unten wird noch vieles über sie gesagt werden. Insbesondere lässt sich vieles aus der Lösung reiner Gleichungen hierher übertragen; darauf werden wir an der jeweiligen Stelle hinzuweisen nicht versäumen. Wir fügen noch ein Beispiel hinzu, dessen Vergleich mit den Regeln den weniger Geübten alles deutlicher machen wird.

Es sei $n=31$, $p=311$, gesucht sind also die Wurzeln der Kongruenz $x^{31}-1\equiv0\;(\mathrm{mod.}\,311)$. Zunächst ist eine primitive Wurzel der Kongruenz $y^{30}-1\equiv\;(\mathrm{mod.}\,31)$ zu suchen, etwa $y\equiv3$. Wir setzen daher $\rho\equiv3$ und teilen zunächst alle primitiven Wurzeln der gegebenen Kongruenz in $5$ Perioden ein, nämlich
\SourceLeaf{220}
\[
\begin{array}{rl}
(5\ast0)\ \dotsfive&(1)+(26)+(25)+(30)+(5)+(6)\\
(5\ast1)\ \dotsfive&(3)+(16)+(13)+(28)+(15)+(18)\\
(5\ast2)\ \dotsfive&(9)+(17)+(8)+(22)+(14)+(23)\\
(5\ast3)\ \dotsfive&(27)+(20)+(24)+(4)+(11)+(7)\\
(5\ast4)\ \dotsfive&(19)+(29)+(10)+(12)+(2)+(21)
\end{array}
\]

Die erforderlichen Rechnungen ergeben für die Summe der Perioden $\equiv-1$, für die ihrer Quadrate $\equiv25$, Kuben $\equiv26$, vierten Potenzen $\equiv249$ und fünften Potenzen $\equiv564$.

Die Perioden sind daher die Wurzeln der Kongruenz
\[
x^5+x^4-12x^3-21x^2+x+5\equiv0
\]

Weiter findet man
\[
\begin{aligned}
(5\ast0)^2&\equiv6+2(5\ast0)+2(5\ast3)+(5\ast4)\\
(5\ast0)^3&\equiv12+15(5\ast0)+4(5\ast1)+3(5\ast2)+6(5\ast3)+6(5\ast4)\\
(5\ast0)^4&\equiv90+60(5\ast0)+28(5\ast1)+26(5\ast2)+49(5\ast3)+38(5\ast4)
\end{aligned}
\]
und daraus durch Elimination
\[
\begin{aligned}
5(5\ast1)&\equiv3(5\ast0)^4-(5\ast0)^3-33(5\ast0)^2-24(5\ast0)+15\\
5(5\ast2)&\equiv-2(5\ast0)^4-(5\ast0)^3+22(5\ast0)^2+31(5\ast0)\\
5(5\ast3)&\equiv(5\ast0)^4-2(5\ast0)^3\\
5(5\ast4)&\equiv-2(5\ast0)^4+4(5\ast0)^3
\end{aligned}
\]

Eine Wurzel der gefundenen Kongruenz ist $\equiv17$. Setzt man also $(5\ast0)\equiv17$, so ist $(5\ast1)\equiv183$, $(5\ast2)\equiv263$, $(5\ast3)\equiv91$, $(5\ast4)\equiv67$.

Nun teile man jede der gefundenen Perioden wiederum in drei, nämlich
\[
\begin{array}{rl}
(5\ast0)&\text{in }(15\ast0),\ (15\ast5),\ (15\ast10)\\
&\text{oder in }(1)+(30),\ (26)+(5),\ (25)+(6)\\
(5\ast1)&\text{in }(15\ast1),\ (15\ast6),\ (15\ast11)\\
&\text{oder in }(3)+(28),\ (16)+(15),\ (13)+(18)\\
&\text{usw.}\qquad\text{usw.}
\end{array}
\]

Die Perioden, in die die folgenden zerlegt sind, seien Wurzeln der jeweils angegebenen Kongruenzen:
\[
\begin{array}{rll}
(5\ast0)&\text{Wurzeln der Kongruenzen sind}&x^3+Ax^2+Bx+C\equiv0\\
(5\ast1)&&x^3+A'x^2+B'x+C'\equiv0\\
(5\ast2)&&x^3+A''x^2+B''x+C''\equiv0\\
&\text{usw.}
\end{array}
\]

\SourceLeaf{221}
Dann ist
\[
\begin{array}{lll}
A\equiv-(5\ast0),&B\equiv(5\ast0)+(5\ast3),&C\equiv-2-(5\ast4)\\
A'\equiv-(5\ast1),&B'\equiv(5\ast1)+(5\ast4),&C'\equiv-2-(5\ast0)\\
\text{usw.}&\text{usw.}&\text{usw.}
\end{array}
\]

\begingroup\setlength{\arraycolsep}{4pt}Daher
\[
\begin{array}{lll}
(15\ast0),\ (15\ast5),\ (15\ast10)&\begin{gathered}\text{sind Wurzeln}\\\text{der Kongruenzen}\end{gathered}&x^3-17x^2+108x-60\equiv0\\
(15\ast1),\ (15\ast6),\ (15\ast11)&&\\
(15\ast2),\ (15\ast7),\ (15\ast12)&&\\
(15\ast3),\ (15\ast8),\ (15\ast13)&&\\
(15\ast4),\ (15\ast9),\ (15\ast14)&&
\end{array}
\]\endgroup

Hier gilt aber
\[
\begin{gathered}
(15\ast0)^3-3(15\ast0)\equiv(15\ast1)\\
(15\ast1)^3-3(15\ast1)\equiv(15\ast2)\\
\text{usw.}
\end{gathered}
\]

Setzt man daher eine der Wurzeln der ersten Kongruenz, $10$, als $(15\ast0)$, so erhält man
\[
\begin{array}{rlrlrl}
(15\ast0)\equiv&10&(15\ast5)\equiv&& (15\ast10)\equiv&\\
(15\ast1)\equiv&37&(15\ast6)\equiv&& (15\ast11)\equiv&\\
(15\ast2)\equiv&-151&(15\ast7)\equiv&& (15\ast12)\equiv&\\
(15\ast3)\equiv&-39&(15\ast8)\equiv&& (15\ast13)\equiv&\\
(15\ast4)\equiv&-112&(15\ast9)\equiv&& (15\ast14)\equiv&
\end{array}
\]

Schließlich nehme man die Terme, aus denen die einzelnen Perioden bestehen. Dann sind
\[
\begin{array}{lll}
(1),\ (30)&\text{Wurzeln der Kongruenzen}&x^2-(15\ast0)x+1\equiv0\\
(3),\ (28)&&x^2-(15\ast1)x+1\equiv0\\
&\text{usw.}
\end{array}
\]

Die Wurzeln der ersten Kongruenz sind $126$ und $195$, die somit primitive Wurzeln der Kongruenz $x^{31}\equiv1$ sind. Aus ihnen lassen sich die übrigen ohne Mühe ableiten.
\BookRule{.28\linewidth}
